\bookchapter{Lampiran A: Alur “Level Nanny” Unboxing Ponsel Baru}

\booksection{“Tiga Langkah” Unboxing dan Pemeriksaan:}

\begin{itemize}
  \item \textbf{Periksa tampilan}: Periksa apakah ada goresan atau benturan pada bodi dan layar.
  \item \textbf{Cocokkan informasi}: Pastikan IMEI/nomor seri pada kotak sesuai dengan yang ada di faktur.
  \item \textbf{Uji fungsi}: Telepon untuk menguji earpiece dan mikrofon, ambil foto untuk memeriksa kamera.
\end{itemize}

\booksection{“Pengingat Penting” Aktivasi dan Pengaturan Awal:}

\begin{itemize}
  \item \textbf{Jaringan}: Hubungkan ke Wi-Fi yang stabil terlebih dahulu.
  \item \textbf{Akun}: Login dengan Apple ID, akun Huawei, akun Xiaomi, dll.
  \item \textbf{Pengaturan dasar}: Aktifkan "Find My Phone", dan atur kata sandi layar kunci yang kuat.
  \item \textbf{Pengaturan eSIM}: Jika ponsel Anda mendukung eSIM (seperti iPhone, beberapa flagship Android), Anda dapat menambahkan paket eSIM di pengaturan tanpa perlu memasukkan kartu SIM fisik. Saat bepergian ke luar negeri, eSIM sangat nyaman karena Anda dapat langsung membeli paket eSIM operator lokal.
  \item \textbf{Migrasi data}: Saat memigrasikan data dari ponsel lama, pastikan kedua ponsel memiliki baterai yang cukup. iPhone gunakan "Quick Start", Android gunakan aplikasi pindah perangkat resmi.
\end{itemize}

\booksection{Memasang Aplikasi dengan “Utamakan Keamanan”:}

\begin{itemize}
  \item \textbf{Utamakan toko resmi}: Apple gunakan App Store, Android gunakan toko aplikasi bawaan ponsel.
  \item \textbf{Jangan berikan izin sembarangan}: Saat muncul permintaan izin saat instalasi, pikirkan baik-baik "Apakah aplikasi ini benar-benar membutuhkannya?"
\end{itemize}
